% This is part of Mes notes de mathématique
% Copyright (c) 2011-2020
%   Laurent Claessens
% See the file fdl-1.3.txt for copying conditions

%+++++++++++++++++++++++++++++++++++++++++++++++++++++++++++++++++++++++++++++++++++++++++++++++++++++++++++++++++++++++++++ 
\section{Groupes}
%+++++++++++++++++++++++++++++++++++++++++++++++++++++++++++++++++++++++++++++++++++++++++++++++++++++++++++++++++++++++++++

%---------------------------------------------------------------------------------------------------------------------------
\subsection{Définition, unicité du neutre}
%---------------------------------------------------------------------------------------------------------------------------

\begin{definition}[Groupe]      \label{DEFooBMUZooLAfbeM}
    Un \defe{groupe}{groupe} est un ensemble \( G\) muni d'une opération interne \( \cdot\colon G\times G\to G\) telle que
    \begin{enumerate}
        \item
            pour tous \( g\), \( h\), \( k\in G\), \( g\cdot(h\cdot k)=(g\cdot h)\cdot k\),
        \item
            il existe un élément \( e\in G\) tel que \( e\cdot g=g\cdot e=g\) pour tout \( g\in G\),
        \item
            pour tout \( g\in G\), il existe un élément \( h\in  G\) tel que \(g\cdot h=h\cdot g=e \).
    \end{enumerate}
\end{definition}

    Notons que nous avons écrit \( g\cdot h\) et non \( \cdot(g,h)\) comme une notation purement fonctionnelle nous l'aurait suggéré. Dans les exemples concrets, selon les cas, le loi de groupe appliquée à \( g\) et \( h\) sera notée tantôt \( g+h\), tantôt \( g\cdot h\) ou, le plus souvent pour un groupe générique, simplement \( gh\).

\begin{lemmaDef}[Unicités]  \label{LEMooECDMooCkWxXf}
    L'inverse et le neutre sont uniques, c'est-à-dire :
    \begin{enumerate}
        \item
            il existe un unique élément \( e\in G\) tel que \( e g=g e=g\) pour tout \( g\in G\),
        \item       \label{ITEMooOIWTooYqmMPP}
            pour tout \( g\in G\), il existe un unique élément \( h\in  G\) tel que \(g h=h g=e \).
    \end{enumerate}
    Le \( e\) ainsi défini est nommé \defe{neutre}{neutre!dans un groupe} de \( G\). Le \( h\) tel que \( g h=h g=e\) est nommé l'\defe{inverse}{inverse!dans un groupe} de \( g\) et est noté \( g^{-1}\).
\end{lemmaDef}

\begin{proof}
    Chaque point séparément.
    \begin{enumerate}
        \item
            Supposons que \( e_1\) et \( e_2\) vérifient la propriété. Nous avons pour tout \( g\in G\) : \( e_1g=ge_1=g\). En particulier pour \( g=e_2\) nous écrivons \( e_1e_2=e_2e_1=e_2\). Mais en partant dans l'autre sens : \( e_2g=ge_2=g\) avec \( g=e_1\) nous avons \( e_2e_1=e_1e_2=e_1\). En égalant ces deux valeurs de \( e_2e_1\) nous avons \( e_1=e_2\).

            Pour la suite de la preuve nous écrivons \( e\) l'unique neutre de \( G\).

        \item
            Supposons que \( k_1\) et \( k_2\) soient deux inverses de \( g\). On considère alors le produit \( k_1 g k_2 \). Puisque \(k_1 g = e \), on a \( k_1 g k_2 = e k_2 = k_2 \); mais, comme \(g k_2 = e \), on a aussi \( k_1 g k_2 = k_1 e = k_1 \). Le produit est donc à la fois égal à \( k_1 \) et à \( k_2 \), et donc \( k_1 = k_2 \).
    \end{enumerate}
\end{proof}

\begin{definition}\label{DEFGroupeAbelien}
    Un groupe \( G\) est \defe{abélien}{abélien!groupe} ou \defe{commutatif}{commutatif!groupe} si pour tous \( g\) et \( h\) dans \( G\), \( gh=hg\).
\end{definition}

%---------------------------------------
\subsection{Puissances dans un groupe}
%---------------------------------------

Voici une première définition de la fonction puissance. Il y en aura d'autres, de plus en plus générales. Voir le thème \ref{THEMEooBSBLooWcaQnR}.
\begin{definition}\label{DEFooGVSFooFVLtNo}
    Si \( G\) est un groupe (multiplicatif), si \( g\in G\) et si \( n\in \eN\), nous définissons \( g^n\) par récurrence :
    \begin{enumerate}
        \item
            \( g^0=e\) (le neutre),
        \item       \label{ITEMooOUIPooGjAgQb}
            \( g^{k+1}=g\cdot g^{k}\).
    \end{enumerate}
    On ajoute de plus que \( g^{-n} = (g^{-1})^n \).
\end{definition}

Le lemme suivant dit que le point \ref{ITEMooOUIPooGjAgQb} de la définition \ref{DEFooGVSFooFVLtNo} aurait pu être écrit \( g^k\cdot g\) au lieu de \( g\cdot g^k\).
\begin{lemma}[\cite{MonCerveau}]        \label{LEMooWPARooYLZlzr}
    Si \( G\) est un groupe, si \( a\in G \) et si \( n\in \eN\) est non-nul, alors
    \begin{equation}
        g^n=g\cdot g^{n-1}=g^{n-1}\cdot g.
    \end{equation}
\end{lemma}
\begin{proof}
Pour \( n = 1 \), c'est trivial; pour \( n = 2 \), c'est facile. Pour des \( n \) plus grands, par récurrence, on suppose que \( g^{n-1} = g^{n-2}\cdot g \), et il vient 
\begin{align}
 g^n = g\cdot g^{n-1} &= g\cdot  g^{n-2}\cdot g  &\text{par hypothèse}\\
 & = g^{n-1}\cdot g &\text{par définition}.
 \end{align}
\end{proof}

%--------------------------------------------------------------------
\subsection{Sous-groupes, sous-groupes engendrés}
%--------------------------------------------------------------------

\begin{definition}
    Soit \( H\) une partie du groupe \( G\). On dit que \( H \) est un \defe{sous-groupe}{sous-groupe}\index{sous-groupe} de \( G \) si l'opération sur \( G\) confère à \( H\) une structure de groupe. En d'autres termes, \( H \) est un sous-groupe de \( G \) si:
\begin{itemize}
    \item le neutre est dans \( H\);
    \item un produit d'éléments de \( H\) reste dans \( H\);
    \item l'inverse d'un élément de \( H\) est dans \( H\).
\end{itemize}
\end{definition}

Un truc sympa, coup-de-poing, pour déterminer si on a affaire à un sous-groupe:

\begin{proposition}    \label{PROP_criteressgroupe}
Une partie non-vide \( H\) du groupe \( G\) est un sous-groupe de \( G \) si et seulement si, pour tous \(x\ ,y \in H\), on a \( xy^{-1} \in H\).
\end{proposition}

\begin{proposition}    \label{PROP_interssgroupes}
L'intersection de sous-groupes est encore un sous-groupe.
\end{proposition}
\begin{proof}
Si \( x \) et \( y \) se trouvent dans plein de sous-groupes, alors le produit \( xy^{-1} \) se trouvera dans ces mêmes sous-groupes, donc à leur intersection.
\end{proof}

\begin{definition}[Sous-groupe engendré]        \label{DefooRDRXooEhVxxu}
    Soit \( A\) une partie du groupe \( G\). Le sous-groupe \defe{engendré}{sous-groupe!engendré}\index{engendré!sous-groupe} par \( A\) est l'intersection de tous les sous-groupes de \( G\) contenant \( A\). Nous notons ce groupe \( \gr(A)\)\nomenclature[R]{\( \gr\)}{groupe engendré}.

    Lorsque \( A \) est fini (disons \( A = \{a_1, \dots, a_n\} \)), on note aussi le sous-groupe engendré \( \langle a_1, \dots, a_n \rangle \).

    En particulier, \( \gr(\emptyset)=\{ e \}\)\footnote{Vu que l'ensemble vide est une partie de tous les ensembles, la définition \( \gr(\emptyset)=\{ e \}\) est ambigüe : c'est l'identité de quel groupe ? C'est le contexte qui dira.}.
\end{definition}

Un sous-groupe engendré n'est jamais vide parce qu'il contient toujours au moins le neutre (parce que c'est un sous-groupe). Si \( G\) est un groupe, le sous-groupe \( \gr(\emptyset)\) lui-même contient \( e\)\footnote{Demandez-vous si il est possible que \( \gr(\emptyset)\) contienne d'autres éléments que \( e\).}.

\begin{lemma}
    Si \( G\) est un groupe et \( A\) une partie de \( G\), alors \( \gr(A)\) est un sous-groupe de \( G\).
\end{lemma}
\begin{proof}
   Par la définition \ref{DefooRDRXooEhVxxu} de \( \gr(A)\), c'est une intersection de sous-groupes, donc un sous-groupe par la proposition \ref{PROP_interssgroupes}.
\end{proof}

Le sous-groupe engendré par \( A \) est le plus petit (pour l'inclusion) groupe de \( G\) contenant \( A\). Plus formellement, nous avons le résultat suivant.
\begin{lemma}     \label{LEM_contenancedussgrengendre}
    Tout sous-groupe de \( G\) contenant \( A\) contient \( \gr(A)\).
\end{lemma}

\begin{proof}
    Si \( H\) est un sous-groupe de \( G\) contenant \( A\), alors \( \gr(A)\) est l'intersection de \( H\) avec tous les autres sous-groupes de \( G\) contenant \( A\). Il contient donc \( \gr(A)\).
\end{proof}

Avant d'aborder une caractérisation explicite d'un sous groupe engendré, voici un peu de notations.
\begin{definition}      \label{DEFooNEVNooJlmJOC}
    Si \( f\colon \{ 0,\ldots, N \}\to G\) est une application vers un groupe additif \( G\), alors nous définissons la notation \( \sum_{i=0}^Nf(i)\) par récurrence de la façon suivante :
    \begin{enumerate}
        \item
            \( \sum_{i=0}^0f(i)=f(0)\),
        \item
            \( \sum_{i=0}^{k}f(i)=\sum_{i=0}^{k-1}f(i)+f(k)\).
    \end{enumerate}
    Dans le cas d'un groupe multplicatif, on utilise la notation \( \prod_{i=0}^N f(i)\) par récurrence également:  \( \prod_{i=0}^0 f(i) = f(0)\), et  \( \prod_{i=0}^{k}f(i)=\prod_{i=0}^{k-1}f(i)\cdot f(k)\).
\end{definition}


\begin{lemma}[\cite{BIBooERNQooTXQPvD}]   \label{LemFUIZooBZTCiy}
    Si \( A\) est une partie du groupe \( G\), alors le sous-groupe \( \gr(A)\) engendré\footnote{Définition~\ref{DefooRDRXooEhVxxu}.} par \( A\) est l'ensemble de tous les produits finis d'éléments de \( A\) et de \( A^{-1}\) (l'identité est le produit à zéro éléments).

    C'est-à-dire que tout élément de \( \gr(A)\) peut être écrit sous la forme
    \begin{equation}\label{EQ_produitGenerateursGroupe}
        \prod_{i=1}^ng_i^{a_i}
    \end{equation}
    où \( a_i\in \eZ\) et \( g\colon\{ 1,\ldots, n \}\to A\) n'est pas spécialement injective : il peut arriver que \( g_i=g_j\).
\end{lemma}
L'écriture \eqref{EQ_produitGenerateursGroupe} a un sens grâce à la définition \ref{DEFooGVSFooFVLtNo}; les \( a_i\) négatifs correspondant aux inverses. Notons que si \( g\in A\), il n'y a pas de garanties que \( g^{-1}\) soit également dans \( A\).

\begin{proof}
    Vu qu'un produit vide est égal à l'identité, le lemme est vrai (un peu trivialement) dans le cas où \( A=\emptyset\). À partir de maintenant nous supposons que \( A\) est non vide.

    Nous nommons \( \gr(A)\) le groupe engendré par \( A\) et par \( H\) l'ensemble
    \begin{equation}
        H=\{ g_1\ldots g_n\tq n \in \eN;  g_i\in A \text{ ou } g_i^{-1} \in A \forall 1 \leq i \leq n\},
    \end{equation}
    constitué des produits d'éléments de \( A \) ou de leurs inverses\footnote{On peut très bien avoir dans un produit des éléments de $A$ et des inverses.}.
    
    Nous commençons par prouver que \( H\) est un groupe.
    \begin{itemize}
        \item Vu que \( A\) est non vide, nous considérons \( a\in A\). Dans ce cas, \( e=aa^{-1}\in H\). Donc \( e\in H\).
        \item L'inverse de \( g_1\ldots g_n\) est \( g_n^{-1}\ldots g_1^{-1}\) qui est également dans \( H\).
        \item Le produit de \( g_1\ldots g_n\) par \( h_1\ldots h_n\) est également dans \( H\)\footnote{Et c'est ici qu'on se rend compte que la décomposition n'est probablement que rarement unique.}.
    \end{itemize}
    
    Voyons à présent que \( H \) n'est autre que le groupe engendré par \( A \); procédons par double inclusion. D'abord, tout groupe contenant \( A\) doit contenir les inverses des éléments de \( A \) et les produits finis d'éléments de \( A \); donc \( H\subset \gr(A)\). Ensuite, \( H\) est un groupe contenant \( A\), donc le lemme \ref{LEM_contenancedussgrengendre} dit que \( \gr(A)\subset H\). Au final, \( H=\gr(A)\), ce qu'il fallait.
\end{proof}

\begin{lemma}       \label{LEMooCFTVooKvmyKN}
    Soit un groupe \( G\) et un sous-groupe \( H=\gr(h_1,\ldots, h_n)\). Si \( \alpha\in G\), alors
    \begin{equation}
        \alpha H\alpha^{-1}=\gr(\alpha h_1\alpha^{-1},\ldots, \alpha h_n\alpha^{-1}).
    \end{equation}
\end{lemma}

\begin{proof}
    Il s'agit d'une conséquence du lemme \ref{LemFUIZooBZTCiy}. Un élément de \( \gr(\alpha h_1\alpha^{-1},\ldots, \alpha h_n\alpha^{-1})\) est un produit d'éléments de \( G\) de la forme \( \alpha h_i\alpha^{-1}\) ou \( (\alpha h_j\alpha^{-1})^{-1}=\alpha h_j^{-1}\alpha^{-1}\). Or nous avons
    \begin{equation}
        \alpha h_i\alpha^{-1}\alpha h_j\alpha^{-1}=\alpha h_ih_j\alpha^{-1}\in \alpha H\alpha^{-1}.
    \end{equation}
    Donc 
    \begin{equation}
        \gr(\alpha h_1\alpha^{-1},\ldots, \alpha h_n\alpha^{-1})\subset \alpha H\alpha^{-1}.
    \end{equation}
    L'inclusion dans l'autre sens est du même tonneau.
\end{proof}

\begin{definition}[Partie génératrice, groupe monogène, générateur, groupe cyclique]  \label{DEFooWMFVooLDqVxR}  \label{DefHFJWooFxkzCF}
    Soit \( G\), un groupe et \( A\) une partie de \( G\). Si \( \gr(A)=G\), alors nous disons que \( A\) est une \defe{partie génératrice}{partie génératrice} du groupe \( G\).

    Un groupe est \defe{monogène}{monogène} s'il a une partie génératrice réduite à un seul élément.
    
    Un tel élément \( a\in G\) est un \defe{générateur}{générateur} de \( G\); tous les éléments de \( G\) s'écrivent alors sous la forme \( a^n\) pour un certain \( n\in\eZ\).
    
    Un groupe fini et monogène est dit \defe{cyclique}{groupe cyclique}.
\end{definition}

\begin{example}
    Soit le groupe \( \big( \eZ/10\eZ,+ \big) \footnote{D'accord, on n'a pas encore défini les groupes quotients, mais ça le fait quand même: pensons qu'on étudie les restes de la division par 10. Comment ça, on n'a pas encore fait de division euclidienne?!} \). L'élément \( [2]_{10}\) n'est pas générateur parce que ses puissances\footnote{Attention aux notations; en général on écrit la loi de groupe de façon multiplicative et on parle des puissances d'un élément, mais ici on écrit la loi de groupe additivement, donc les «puissances» sont en réalité les multiples.} sont
    \begin{equation}
        \gr([2]_{10})=\{ [2]_{10},[4]_{10},[6]_{10},[8]_{10},[0]_{10} \}.
    \end{equation}
    Par contre l'élément \( [3]_{10}\) est générateur : ses puissances sont dans l'ordre
    \begin{equation}
        [3]_{10}, [6]_{10}, [9]_{10}, [2]_{10}, [5]_{10}, [8]_{10},[1]_{10},[4]_{10},[7]_{10},[0]_{10}.
    \end{equation}
\end{example}

Un exemple presque identique, mais un peu masqué sera l'exemple \ref{EXooOXAAooZMdDfP}.

\begin{proposition}[\cite{PDFpersoWanadoo}]
  \begin{enumerate}
  \item
    Un groupe monogène est abélien.
  \item
    Un groupe monogène qui n'est pas fini, est isomorphe à \( \eZ \), le groupe additif des entiers.
  \item
    Un groupe cyclique est isomorphe au groupe additif \( \eZ / n\eZ \) avec \( n \) le nombre d'éléments du groupe.
  \end{enumerate}
\end{proposition}

%---------------------------------------------------------------------------------------------------------------------------
\subsection{Morphismes de groupes}
%---------------------------------------------------------------------------------------------------------------------------
\begin{lemma}       \label{LEMooBIBFooBHxFYC}
    Si \( G\) est un groupe et si \( h\in G\), alors les applications
    \begin{equation}
        \begin{aligned}
            L_h\colon G&\to G \\
            g&\mapsto hg 
        \end{aligned}
    \end{equation}
    et
    \begin{equation}
        \begin{aligned}
            R_h\colon G&\to G \\
            g&\mapsto gh 
        \end{aligned}
    \end{equation}
    sont des bijections (mais ne sont pas des morphismes\footnote{La définition suit ce lemme!} quand \(h \neq e \)).
\end{lemma}

\begin{proof}
    Ce ne sont, en effet, pas des morphismes, puisque \(e \) n'est pas envoyé sur \(e \), sauf cas bien particulier\dots\ Montrons à présent que ce sont tout de même des bijections.

    D'abord si \( L_h(g_1)=L_h(g_2)\), alors \( hg_1=hg_2\) et en multipliant à gauche par \( h^{-1}\) nous avons \( g_1=g_2\); donc \( L_h\) est injective. Ensuite \( L_h\) est surjective parce que si \( g\in G\), alors \( g=L_h(h^{-1} g)\).

    Pour l'application \( R_h\), la preuve est une simple adaptation.
\end{proof}

\begin{definition}
    Soit \(G,\ H\) deux groupes, et \(f \) une application de \(G \) dans \(H \). Si \(f \) respecte les structures de groupe, alors \(f \) est appelé un \defe{morphisme de groupes}{morphisme!de groupes} (ou un \emph{homomorphisme de groupes}). Respecter les structures de groupe signifie que:
    \begin{enumerate}
    \item \( f(e_G) = e_H \): le neutre de \( G \) est envoyé sur le neutre de \( H \);
    \item pour tous \(g, g' \in G,\ f(gg') = f(g)f(g')\): l'image d'un produit (dans \(G \)) est le produit (dans \( H \)) des images;
    \item pour tout \(g \in G,\ f(g^{-1}) = \left(f(g)\right)^{-1}\): l'image de l'inverse d'un élément est l'inverse de l'image de cet élément.
    \end{enumerate}
\end{definition}


\begin{definition}      \label{DEFooWBIYooGNRYOp}
    Soient un groupe \( G\), un ensemble $X$ et une application \( f\colon X\to G\). Le \defe{noyau}{noyau!application vers un groupe} de \( f\) est la partie
    \begin{equation}
        \ker(f)=\{ x\in X\tq f(x)=e \}
    \end{equation}
    où \( e\) est l'unité de \( G\).
    
    L'\defe{image}{image!d'un morphisme} de \(f \) est l'ensemble 
    \begin{equation}
        \Image(f)=f(X) = \{ f(x) \tq x\in X \}.
    \end{equation}
\end{definition}

\begin{remark}
  Quand \( f \) est un morphisme de groupes, son noyau et son image sont des sous-groupes de leurs groupes respectifs. Nous le prouverons dans le cadre du premier théorème d'isomorphisme \ref{ThoPremierthoisomo}.
\end{remark}

\begin{definition}\label{DEFGroupesTypesdeMorphismes}
Un morphisme injectif s'appelle un \emph{monomorphisme}.

Un morphisme surjectif s'appelle un \emph{épimorphisme}.

Un morphisme bijectif s'appelle un \emph{isomorphisme}.

Les morphismes d'un groupe dans lui-même sont appelés des \emph{endomorphismes}; l'ensemble des endomorphismes d'un groupe est noté \( \End(G) \).

Quand ils sont bijectifs, on dit que ce sont des \emph{automorphismes}.  L'ensemble des automorphismes d'un groupe est noté \( \Aut(G) \).
\end{definition}

\begin{lemma}
Un morphisme est injectif si et seulement si son noyau est réduit au neutre.

Un morphisme est surjectif si et seulement si son image est le groupe d'arrivée.
\end{lemma}

\begin{definition}      \label{DEFooUXXTooCCLmQe}
    Un sous-groupe \( H\) de \( G\) est un sous-groupe \defe{caractéristique}{sous-groupe!caractéristique}\index{caractéristique!sous-groupe} s'il est stable pour tous les automorphismes\footnote{Définition \ref{DEFGroupesTypesdeMorphismes}.} de \(G \); en d'autres termes, un sous-groupe caractéristique \(H \) de \(G \) vérifie \( \alpha(H)=H\) pour tout \( \alpha\in \Aut(G)\).
\end{definition}


%---------------------------------------------------------------------------------------------------------------------------
\subsection{Classes de conjugaison}
%---------------------------------------------------------------------------------------------------------------------------

\begin{definition}
    Soit un groupe \( G\) et un élément \( g\in G\). La \defe{classe de conjugaison}{classe!de conjugaison} de \( g\) est la partie
    \begin{equation}
        C_g=\{ kgk^{-1}\tq k\in G \}.
    \end{equation}
\end{definition}

\begin{lemma}       \label{LEMooQYBJooYwMwGM}
    Un groupe est abélien si et seulement si ses classes de conjugaison sont des singletons.
\end{lemma}

\begin{proof}
    Supposons que \( G\) soit abélien. Alors
    \begin{equation}
        C_g=\{ kgk^{-1}\tq k\in G \}=\{ g \}.
    \end{equation}
    Donc les classes de conjugaison sont des singletons.

    Dans l'autre sens, si les classes sont des singletons, on a \( kgk^{-1}=g\) pour tous \( k,g\in G\). Cela signifie immédiatement que \( G\) est abélien.
\end{proof}


%--------------------------------------
\subsection{Sous-groupe normal, groupes quotients}
%--------------------------------------

\begin{definition}\label{DEFooNIIMooFkZgvX}
    Un sous-groupe \( N\) de \( G\) est \defe{normal}{normal!sous-groupe} ou \defe{distingué}{distingué!sous-groupe} (dans \( G \) si pour tout \( g\in G\) et pour tout \( n\in N\), \( gng^{-1}\in N\). Autrement dit lorsque \( gNg^{-1}\subset N\).

    Lorsque \( N\) est normal dans \( G\) il est parfois noté \( N\normal G\)\nomenclature[]{\(N \normal G\)}{Le sous-groupe \( N\) est normal dans \( G\)}.
\end{definition}

\begin{example}
  Les noyaux de morphismes de groupes sont des sous-groupes normaux. Nous le prouverons dans quelques instants, avec le théorème du premier isomorphisme \ref{ThoPremierthoisomo}.
\end{example}

\begin{proposition}\label{propGroupeNormal}
    Soit \( N\) un sous-groupe de \( G\). Les propriétés suivantes sont équivalentes :
    \begin{enumerate}
        \item
            \( gNg^{-1}= N\) pour tout \( g\in G\),
        \item
            \( gN=Ng\) pour tout \( g\in G\),
        \item
            pour tout \( h \in N \), la classe de conjugaison \( C_h \) est dans \( N \),
        \item
            \( N\) est une union de classes de conjugaison de \( G\),
        \item
            \( N\) est normal\footnote{Définition \ref{DEFooNIIMooFkZgvX}.}.
        \item
            \( N \) est normal dans tout sous-groupe qui le contient.
    \end{enumerate}
\end{proposition}

\begin{lemma}       \label{LEMooFNVRooRCkjLc}
    Lorsque \( N\) est normal dans \( G\), alors la définition
    \begin{equation}        \label{EQooEUESooSeUWHK}
        [a]\cdot[b]=[ab]
    \end{equation}
    définit une loi de groupe sur l'ensemble \( G/N\).
\end{lemma}

\begin{proof}
    Le neutre est \( [e]\) et l'associativité ne pose pas plus de problèmes que l'existence d'un inverse. Le point à vérifier est que la formule \eqref{EQooEUESooSeUWHK} est une bonne définition : \( [ah]\cdot [bh']=[ab]\) pour tout \( h,h'\in N\). Nous avons :
    \begin{equation}
        [ah]\cdot [ah']=[ahah']=[ahb].
    \end{equation}
    Pour montrer que cela est \( [ab]\), l'astuce est d'introduire \( bb^{-1}\) à côté du \( a\) :
    \begin{equation}
        [ahb]=[abb^{-1}hb]=[ab]
    \end{equation}
    parce que \( b^{-1} hb\in N\) du fait que \( N\) soit normal dans \( G\).
\end{proof}

\begin{example}[\cite{ooTVPEooGOLEom}]      \label{EXooFNIKooHxePSs}
    Il ne faudrait pas croire que le groupe quotient \( G/N\) est forcément un sous-groupe de \( G\): déjà, ils s'agit d'objets différents. Il ne faudrait pas croire non plus que \( G/N\) soit isomorphe à un sous-groupe de \( G \). Par exemple, le quotient \( \eZ/2\eZ\) est isomorphe à l'ensemble \( \{ 0,1 \}\), muni de l'addition telle que \( 1+1=0\); cette égalité est évidemment fausse dans \( \eZ\).

    En revanche, il est bien vrai que l'application
    \begin{equation}
        \begin{aligned}
            \pi \colon G&\to G/N \\
            g&\mapsto [g]
        \end{aligned}
    \end{equation}
    est un morphisme de groupes; on l'appelle projection canonique.
\end{example}

\begin{proposition}[Caractérisation des sous-groupes d'un groupe quotient -- \cite{ooTVPEooGOLEom}]
Soit \( G \) un groupe, \( N \) un sous-groupe normal de \( G \), et  \( \pi\colon G\to G/N\) la projection cannique (de l'exemple précédent).
  \begin{enumerate}
  \item
    \( \pi \)  est un homomorphisme surjectif de noyau~\( N\).
  \item
      Les sous-groupes de \( G \) contenant \( N \) sont en bijection, par  \( \pi \), avec les sous-groupes de \( G/N \).
  \item
      Les sous-groupes normaux de \( G \) contenant \( H \) sont en bijection, par  \( \pi \), avec les sous-groupes normaux de \( G/H \).
  \end{enumerate}
\end{proposition}


\begin{definition}[Groupe simple]   \label{DefGroupeSimple}
    Un groupe est dit \defe{simple}{simple!groupe} si il est non trivial et si les seuls sous-groupes normaux qu'il admet sont lui-même et le sous-groupe réduit à l'identité.
\end{definition}

\begin{lemma}[\cite{PDFpersoWanadoo}]\label{LemHUkMxp}
    Si \( H\) et \( K\) sont normaux dans le groupe \( G\) et si \( H\cap K=\{ e \}\) alors \( HK\simeq H\times K\).
\end{lemma}

\begin{proposition} \label{PropSRMJooIDPBoW}
    Soit deux groupes \( G\) et \( H \) . Nous considérons un sous-groupe normal \( N\) de \( G\) ainsi qu'un homomorphisme \( \psi\colon G\to H\). Alors
    \begin{enumerate}
        \item
            \( \psi(N)\) est normal dans \( \psi(G)\)
        \item
            Si \( G/N\) est abélien alors \( \psi(G)/\psi(N)\) est abélien.
    \end{enumerate}
\end{proposition}

\begin{proof}
    Soient \( k\in N\) et \( g\in G\). Alors \( \psi(g)\psi(k)\psi(g)^{-1}=\psi(gkg^{-1})\in\psi(N)\). Donc \( \psi(N)\) est normal dans \( \psi(G)\).

    Pour la seconde partie nous notons \( [\ldots]\) les classes par rapport à \( H\) et \( \overline{ \vphantom{g}\ldots }\) celles par rapport à \( G\). Nous avons, pour tous \( g_1\ , g_2 \in G\):
    \begin{subequations}
        \begin{align}
            [\psi(g_1)][\psi(g_2)]&=\big[ \psi(g_1)\psi(g_2) \big]\\
            &=\big[ \psi(g_1g_2) \big]\\
            &=\{ \psi(g_1g_2)\psi(k)\tq k\in N \}\\
            &=\{ \psi(g_1g_2k)\tq k\in N \}\\
            &=\psi\Big(  \{ g_1g_2k\tq k\in N \}  \Big) \\
            &=\psi\big( \overline{ g_1g_2 } \big)\\
            &=\psi(\overline{ g_2g_1 })\\
            &=\text{refaire à l'envers}\\
            &=[\psi(g_2)][\psi(g_1)].
        \end{align}
    \end{subequations}
    Par conséquent \( \psi(G)/\psi(N)\) est abélien.
\end{proof}


%-----------------------------------------------
\subsection{Théorèmes d'isomorphismes}
%-----------------------------------------------

\begin{theorem}[Premier théorème d'isomorphisme]        \label{ThoPremierthoisomo}
    Soit \( \theta\colon G\to H\) un homomorphisme de groupe. Alors
    \begin{enumerate}
        \item
            \( \Kernel\theta\) est normal dans \( G\),
        \item
            \( \Image \theta\) est un sous-groupe de \( H\)
        \item   \label{ItemWLCLdk}
            nous avons un isomorphisme naturel
            \begin{equation}
                G/\Kernel\theta\simeq \Image\theta
            \end{equation}
    \end{enumerate}
\end{theorem}
\index{théorème!isomorphisme!premier!pour les groupes}

\begin{proof}
    Point par point.
    \begin{enumerate}
        \item
            Le fait que  \( \Kernel\theta\) soit un sous-groupe de \( G\) est clair; montrons qu'il est normal. Si \( g \in G \) et \( u \in \Kernel\theta\), alors \(\theta (g^{-1} u g) = \theta(g^{-1})\theta(u)\theta(g) = \bigl(\theta(g)\bigr)^{-1}\theta(g) = 1_H \), et donc \( g^{-1} u g \in \Kernel\theta\).
        \item
            Il suffit de remarquer que si \( h = \theta(g) \) et \( h' = \theta(g') \), alors \( h^{-1} h' = \theta(g^{-1} g') \).
        \item
            Si \( [g]\) représente la classe de \( g\) dans \( G/\Kernel\theta\), l'isomorphisme est donné par \( \varphi[g]=\theta(g)\).
    \end{enumerate}
\end{proof}

\ifbool{isGiulietta}{
Ce premier théorème d'isomorphismes permet entre autres de prouver que $SO(3)=\SU(2)/\eZ_2$, voir la proposition \ref{PROPooDKPTooBnLflt}.
}{}
\begin{theorem}[Deuxième théorème d'isomorphisme]
    Soient \( H\) et \( N\) deux sous-groupes de \( G\) et supposons que \( N\) soit normal. Alors
    \begin{enumerate}
        \item
            \( NH=HN\) est un sous-groupe.
        \item
            Le groupe \( N\) est normal dans \( NH\).
        \item
            Le groupe \( N\cap H\) est normal dans \( H\).

        \item\label{ItemjRPajc}
            Nous avons l'isomorphisme
            \begin{equation}
                \frac{ HN }{ N }\simeq\frac{ H }{ H\cap N }.
            \end{equation}
        \item   \label{ItembgDQEN}
            Tout les points précédents (et en particulier l'isomorphisme du point~\ref{ItemjRPajc}) sont encore valables si \( N\) n'est pas normal mais si seulement \( H\) normalise \( N\), c'est-à-dire si \( hNh^{-1}\subset N\) pour tout \( h\in H\).
    \end{enumerate}
\end{theorem}
\index{théorème!isomorphisme!second}

\begin{proof}
    Point par point.
    \begin{enumerate}
        \item
            Il est clair que \( 1_G \in NH \). Soient $nh$ et $n'h'$ deux éléments de \( NH \); alors en tenant compte du fait que \( N\) est normal,
            \begin{equation}
                nhn'h'=n\underbrace{hn'h^{-1}}_{\in N}hh'\in NH.
            \end{equation}
            Cela prouve que \( NH\) est un groupe.

            De la même façon, nous prouvons que \( HN\) est un groupe par
            \begin{equation}
                hnh'n'=hh'\underbrace{h'^{-1}nh'}_{\in N}n'\in HN
            \end{equation}

            Nous devons encore prouver que \( HN=NH\). Pour cela, \( nh \in HN \), car \( nh = hh^{-1}nh \), les trois derniers facteurs formant un  élément de \( N \) par normalité; de même \( hn \in NH \), montrant que \( NH = HN \). Enfin, comme \( (nh)^{-1} = h^{-1} n^{-1} \), les inverses de \( NH \) sont dans \( HN = NH \).
        \item
            \( N\) est normal dans \( G \), a fortiori dans l'un de ses sous-groupes. Plus en détail (cela sert pour le point \ref{ItembgDQEN}): si \( k \in N \) et \( nh \in NH \), alors on a \( (nh)^{-1}knh = h^{-1} (n^{-1}kn) h \in N \) car \(n^{-1}kn\) est dans \( N \) et  \( hNh^{-1}\subset N\) (c'est un sous-groupe normal).
        \item
            Il suffit de voir que, si \( h \in H \) et \( n \in N \cap H \), alors \( hnh^{-1} \in N \cap H \). Or, \( hnh^{-1} \in H \) puisque \( H\) est un sous-groupe; et \( hnh^{-1} \in N \) car \( N \) est un sous-groupe normal.
        \item
            Il faut d'abord remarquer que \( H\) et \( N\) étant des groupes et le produit \( NH\) étant un groupe, nous avons \( NH=HN\). Soit le morphisme injectif
            \begin{equation}
                \begin{aligned}
                    j\colon H&\to HN \\
                    h&\mapsto h
                \end{aligned}
            \end{equation}
            et la surjection canonique
            \begin{equation}
                \sigma\colon HN\to HN/N
            \end{equation}
            Nous considérons ensuite l'application composée
            \begin{equation}
                \begin{aligned}
                    f\colon H&\to HN/N \\
                    h&\mapsto hN.
                \end{aligned}
            \end{equation}

            \begin{subproof}
                \item[\( f\) est surjective]

                    L'application \( f\) est surjective parce que l'élément \( hnN\in HN/N\) est l'image de \( h\), étant donné que \( hnN=hN\).

                \item[\( \Kernel(f)=H\cap N\)]

            Si \( a\in H\cap N\), nous avons \( f(a) =aN = N\), et donc \( H\cap N\subset \Kernel f\). D'autre part, si \( h\in H\) vérifie \( h\in\Kernel f\), alors \( f(h)=hN=N\), ce qui est uniquement possible si \( h\in N\).

            \end{subproof}
            Le premier théorème d'isomorphisme implique alors que \( H/\Kernel f\simeq \Image f\), c'est-à-dire
            \begin{equation}
                H/N\cap H\simeq HN/N.
            \end{equation}
       \item
          Les preuves précédentes fonctionnent aussi sous l'hypothèse \( N \normal H \) plutôt que \( N \normal G \): en effet, à chaque fois que l'on a invoqué la normalité, on avait en réalité besoin de justifier \( h^{-1}nh \in N \) pour tout \( n \in N \) et tout \( h \in H \).
    \end{enumerate}
\end{proof}

\begin{theorem}[Troisième théorème d'isomorphisme]  \label{ThoezgBep}
    Soient \( N\) et \( M\) deux sous-groupes normaux de \( G\) avec \( M\subset N\). Alors \( N/M\) est normal dans \( G/M\) et
    \begin{equation}
        (G/M)/(N/M)\simeq G/N.
    \end{equation}
\end{theorem}
\index{théorème!isomorphisme!troisième}

\begin{proof}
    Afin de montrer que \( N/M\) est normal dans \( G/M\), nous considérons \( g\in G\), \( nM\in N/M\) et nous calculons
    \begin{equation}
        gnMg^{-1}=gng^{-1}\underbrace{gMg^{-1}}_{=M}=\underbrace{gng^{-1}}_{\in N}M\in N/M.
    \end{equation}
    Pour prouver l'isomorphisme nous considérons le morphisme
    \begin{equation}
        \begin{aligned}
            \varphi\colon G/M&\to G/N \\
            gM&\mapsto gN.
        \end{aligned}
    \end{equation}
    C'est surjectif et le noyau est \( N/M\) parce que \( \varphi(gM)=N\) uniquement si \( g\in N\). Nous pouvons appliquer le premier théorème d'isomorphisme à \( \varphi\) en écrivant
    \begin{equation}
        (G/M)/\Kernel \varphi\simeq\Image \varphi,
    \end{equation}
    c'est-à-dire
    \begin{equation}
        (G/M)/(N/M)\simeq G/N.
    \end{equation}
\end{proof}

%--------------------------------------------------------------
\subsection{Indice d'un sous-groupe et ordre des éléments}
%--------------------------------------------------------------
\begin{definition}[Ordre d'un groupe et d'un élément]       \label{DEFooKWBCooMlmpCP}
    Ce sont deux choses différentes.
    \begin{enumerate}
        \item

    Si \( G\) est un groupe, l'\defe{ordre}{ordre!d'un groupe} est la cardinalité de \( G\) et est noté \( | G |\).
\item

    L'\defe{ordre}{ordre!élément} d'un élément \( g\) de \( G\) est le naturel
    \begin{equation}
        \min\{ n\in\eN\tq g^n=e \},
    \end{equation}
    s'il existe; dans le cas contraire, nous disons que l'ordre de \( g\) est infini.
    \end{enumerate}
\end{definition}

Nous verrons que le corollaire~\ref{CorpZItFX} au théorème de Lagrange dira que l'ordre d'un élément divise l'ordre du groupe.



\begin{definition}      \label{DEFooMPIAooIeZNaR}
    Si \( H\) est un sous-groupe d'un groupe fini l'\defe{indice}{indice} de \( H\) dans \( G\) est le nombre \( | G |/| H |\), souvent noté \( | G:H |\).
\end{definition}

Le théorème de Lagrange dira en particulier que l'indice est toujours un nombre entier. C'est à ne pas confondre avec le degré d'une extension de corps (définition~\ref{DefUYiyieu}).


\begin{theorem}[Théorème de Lagrange]    \label{ThoLagrange}
    Soit \( H\) un sous-groupe du groupe fini \( G\).  Alors
    \begin{enumerate}
        \item   \label{ITEMooDPKSooNpOusd}
    l'ordre de \( H\) divise l'ordre de \( G\).
\item
    Les trois nombres suivants sont égaux :
    \begin{itemize}
        \item
            le nombre de classes de \( H\) à gauche,
        \item
            le nombre de classes de \( H\) à droite,
        \item
            l'indice de \( H\) dans \( G\).
    \end{itemize}
    \end{enumerate}
    En particulier si \( H\) est distingué dans \( G\) nous avons
    \begin{equation}
        | G/H |=\frac{ | G | }{ | H | }.
    \end{equation}
\end{theorem}
\index{théorème!Lagrange}
\index{Lagrange!théorème}

\begin{proof}
    Nous commençons par montrer que les classes de \( H\) ont toutes le même nombre d'éléments que \( H\). En effet pour chaque \( g\in G\) nous avons la bijection
    \begin{equation}
        \begin{aligned}
            \varphi\colon H&\to gH \\
            h&\mapsto gh.
        \end{aligned}
    \end{equation}
    L'injectivité de \( \varphi\) est le fait que \( gh=gh'\) implique \( h=h'\). La surjectivité est par définition de la classe.

    Les classes à gauche formant une partition de \( G\), le cardinal de \( G\) est le produit de la taille des classes par le nombre de classes :
    \begin{equation}
        | G |=| H |\cdot\text{nombre de classes}.
    \end{equation}
    En particulier nous voyons que \( | H |\) divise \( | G |\).

    La dernière formule exprime simplement que \( G/H\) est par définition le nombre de classes de \( H\) à gauche (ou à droite) dans \( G\).
\end{proof}

\begin{corollary}       \label{CorpZItFX}
    L'ordre d'un élément d'un groupe fini divise l'ordre du groupe. En particulier dans un groupe d'ordre \( n\) tous les éléments vérifient \( q^n=e\).
\end{corollary}

\begin{proof}
    Soit \( G\) un groupe fini et considérons, à \( g \in G \) fixé, le sous-groupe
    \begin{equation}
        H=\{ g^k\tq k\in\eN \}.
    \end{equation}
    Par le théorème de Lagrange~\ref{ThoLagrange}, l'ordre de \( H\) divise \( | G |\), mais l'ordre de \( H\) est le plus petit \( k\) tel que \( g^k=e\), c'est-à-dire l'ordre de \( g\).
\end{proof}

D'autres résultats à propos d'ordres et d'indices de groupes finis dans la proposition \ref{PROPooVWVIooQzuAlA} et le lemme \ref{LemqAUBYn}. En particulier le théorème de Cauchy \ref{THOooSUWKooICbzqM} qui dit si \( p\) divise l'ordre du groupe \( G\), alors \( G\) contient au moins un élément d'ordre \( p\).


Il y a plein d'autres choses à dire sur les groupes. On y reviendra entre autres au chapitre~\ref{CHAPGroupes}.
